\chapter*{Abstract}
\addcontentsline{toc}{chapter}{Abstract}

Rapid and largely unregulated industrialization in Bangladesh has created severe, spatially concentrated air pollution hotspots inside factories and workshops -- welding bays, boiler rooms, generator rooms, chemical storage areas, and factory entrances -- where workers are exposed to hazardous concentrations of particulate matter (PM1.0, PM2.5, PM10), carbon monoxide (CO), nitrogen dioxide (NO$_2$), volatile organic compounds (VOC), ammonia (NH$_3$), hydrogen sulphide (H$_2$S), and smoke for prolonged periods. Conventional air quality management strategies in this context are almost exclusively city-scale or plant-scale, relying on either large centralized purification plants or coarse regulatory monitoring, neither of which protects a worker standing directly beside a hazardous point source. This thesis proposes \textbf{NirmalNode}, a Green Internet of Things (IoT) and Adaptive Artificial Intelligence (AI)-assisted system that performs \emph{local, hotspot-level} air purification rather than area-wide purification, activating a compact filtration unit only at the specific micro-location where a worker is exposed, and only when it is needed.

NirmalNode is built around a low-cost ESP32 microcontroller node instrumented with a PMS5003 particulate sensor, MiCS-4514 and MQ-series gas sensors, and an optional BME280 environmental sensor, which continuously streams pollutant, temperature, humidity, and pressure readings, tagged with time, day, and location identifiers, to a cloud database over Wi-Fi. A lightweight, edge-deployable Adaptive AI model -- trained initially on historical Bangladesh air-quality data and subsequently updated using \emph{incremental (online) learning} on live sensor streams -- forecasts pollutant concentrations one to two hours ahead of the current time. When the forecast crosses a hazard threshold, the system pre-emptively activates a DC fan that draws air through a cyclone separator, a HEPA filter, and an activated-carbon filter stage before releasing purified air back into the worker's immediate breathing zone, so that purification begins \emph{before} pollution becomes dangerous rather than reacting after exposure has already occurred.

The central methodological contribution of this work is the use of incremental/online learning models (including SGD Regressor, Passive-Aggressive Regressor, Hoeffding Tree, and Adaptive Random Forest, implemented using the River and scikit-learn libraries) in place of periodic full retraining. This allows a single deployed unit to gradually adapt its internal pollution model to a new site's characteristics -- for example, when a unit is relocated from an industrial cluster in Gazipur to one in Narayanganj -- without requiring an engineer to manually retrain and redeploy firmware, and while consuming a fraction of the computational and energy budget of a conventional deep-learning retraining cycle. This area-specific adaptivity, combined with the deliberately lightweight (``Green AI'') model choice and the low-power ESP32-centred hardware (``Green IoT''), is intended to make the system economically and technically feasible for widespread deployment across the many small and medium industrial units that characterize Bangladesh's manufacturing sector.

A bench-scale prototype of NirmalNode was designed, instrumented, and evaluated in a controlled laboratory setting and, where feasible, at a representative industrial site, with performance assessed along three axes: (i) predictive accuracy of the adaptive pollution forecasting model, measured using Mean Absolute Error (MAE), Root Mean Squared Error (RMSE), Precision, Recall, F1-score, and coefficient of determination (R$^2$); (ii) purification efficiency of the cyclone--HEPA--activated-carbon filter train, measured as percentage reduction in PM and gaseous pollutant concentration; and (iii) system-level power consumption and estimated deployment cost, to establish the practical viability of the design for resource-constrained industrial settings.

The findings and design decisions of this work are discussed against the backdrop of the existing literature on industrial air monitoring, indoor and portable air purification, and AI/IoT-based air quality prediction reviewed in Chapter 2, and the identified research gap -- the absence of a low-cost, adaptive, hotspot-targeted purification system suited to the Bangladeshi industrial context -- is used to motivate the specific contributions of NirmalNode. The thesis concludes with a discussion of the limitations of the present prototype and outlines a future work agenda spanning multi-node factory-wide deployment, federated learning across nodes, further edge-AI optimization, and a solar-powered variant for use in areas with unreliable grid power.

\vspace{1cm}
\noindent \textbf{Keywords:} Green IoT; Green AI; Adaptive AI; Incremental Learning; Air Pollution Prediction; Local Air Purification; ESP32; HEPA Filter; Activated Carbon; Cyclone Separator; Industrial Hotspot Monitoring; Edge--Cloud Architecture; Bangladesh.
